\newpage
\section{Master Configuration Guide}

To deploy a robust, universal AI agent toolkit, you can combine the best servers from both the Node (npx) and Python (uvx) ecosystems into a single master configuration block. 

When functionalities overlap (such as the \lstinline[style=inline]|fetch| capability), the Node implementation (\lstinline[style=inline]|@modelcontextprotocol/server-fetch|) is typically chosen due to its widespread adoption and canonical support, while \lstinline[style=inline]|uvx| is preferred for native local system tasks (like file I/O and Git history).

\subsection{The Universal JSON Profile}
This configuration aggregates repository control, system auditing, web discovery, and reasoning. Replace the placeholder environment variables with your actual API keys and local paths.

\begin{lstlisting}[style=json]
{
  "mcpServers": {
    "github": {
      "command": "npx",
      "args": ["-y", "@modelcontextprotocol/server-github"],
      "env": { "GITHUB_PERSONAL_ACCESS_TOKEN": "YOUR_TOKEN" }
    },
    "git-history": {
      "command": "uvx",
      "args": ["mcp-server-git"]
    },
    "local-filesystem": {
      "command": "uvx",
      "args": ["mcp-server-filesystem", "/"]
    },
    "fetch-markdown": {
      "command": "npx",
      "args": ["-y", "@modelcontextprotocol/server-fetch"]
    },
    "brave-search": {
      "command": "npx",
      "args": ["-y", "@modelcontextprotocol/server-brave-search"],
      "env": { "BRAVE_API_KEY": "YOUR_KEY" }
    },
    "postgres-db": {
      "command": "npx",
      "args": ["-y", "@modelcontextprotocol/server-postgres"],
      "env": {
        "PGHOST": "localhost",
        "PGPORT": "5432",
        "PGDATABASE": "your_db",
        "PGUSER": "your_user",
        "PGPASSWORD": "your_password"
      }
    },
    "memory": {
      "command": "npx",
      "args": ["-y", "@modelcontextprotocol/server-memory"]
    },
    "sequential-thinking": {
      "command": "npx",
      "args": ["-y", "@modelcontextprotocol/server-sequential-thinking"]
    }
  }
}
\end{lstlisting}

\subsection{Deployment Instructions}
To apply this universal configuration, copy the JSON block above and place it in the respective configuration file for your AI client.

\begin{itemize}
    \item \textbf{Claude Desktop}: 
    \begin{itemize}
        \item \textit{macOS/Linux}: \lstinline[style=inline]|~/.config/Claude/claude_desktop_config.json|
        \item \textit{Windows}: \lstinline[style=inline]|%APPDATA%\Claude\claude_desktop_config.json|
    \end{itemize}
    
    \item \textbf{Gemini CLI}: 
    \begin{itemize}
        \item Place the JSON inside \lstinline[style=inline]|~/.gemini/settings.json| or a local \lstinline[style=inline]|.gemini/settings.json| at your project root.
        \item Verify successful bindings by running \lstinline[style=inline]|/mcp list| inside the prompt.
    \end{itemize}
    
    \item \textbf{Kiro IDE}:
    \begin{itemize}
        \item Map the configuration to \lstinline[style=inline]|~/.config/Kiro/mcp_config.json|.
    \end{itemize}
    
    \item \textbf{Cursor IDE}:
    \begin{itemize}
        \item While Cursor manages configurations via a UI, you can manually input these servers by navigating to \textbf{Settings -> Features -> MCP} and adding each entry using the ``command'' type matching the JSON definitions.
    \end{itemize}
\end{itemize}
